\section{Conclusiones}

La realización de esta práctica permitió caracterizar el funcionamiento de un transformador monofásico mediante ensayos de rutina y determinar los parámetros principales de su circuito equivalente. Se comprobó que el devanado de alta tensión presenta una resistencia mayor que el de baja tensión y esto puede deberse a que utiliza conductores más finos y de mayor longitud. Asimismo, se obtuvo una relación de transformación global de $a=2,29$ y se confirmó que el transformador posee polaridad aditiva, ya que $E_x>E_p$. En el ensayo de vacío se determinaron los parámetros de la rama de excitación, $R_c=240\,\Omega$ y $X_m=80\,\Omega$; además, el bajo factor de potencia ($\cos\phi_0\approx0,313$) evidenció el predominio de la corriente magnetizante y la importancia de las pérdidas en el núcleo. Por otra parte, en el ensayo de cortocircuito se obtuvieron $Z_{cc}=1,38\,\Omega$, $R_{cc}=1,25\,\Omega$ y $X_{cc}=0,58\,\Omega$. Como $R_{cc}>X_{cc}$, las pérdidas resistivas tienen mayor influencia que los efectos reactivos. En conjunto, los ensayos permitieron construir el circuito equivalente completo del transformador y obtener información suficiente para estimar su comportamiento, regulación de tensión y pérdidas bajo diferentes condiciones de carga. Los resultados son coherentes con un transformador monofásico de baja potencia y validan la metodología experimental empleada.